% =====================================================================
%  quy gán nguồn cho rag hai tầng — slide
%  biên dịch:  build/tools/tectonic.exe slides.tex
% =====================================================================
\documentclass[aspectratio=169,10pt]{beamer}

\usepackage{fontspec}
\setmainfont{texgyrepagella}[
  Extension      = .otf,
  UprightFont    = *-regular,
  BoldFont       = *-bold,
  ItalicFont     = *-italic,
  BoldItalicFont = *-bolditalic ]
\setmonofont{lmmono10-regular}[Extension=.otf, Scale=0.82]

\usepackage{tcolorbox}
\usepackage{fancyvrb}
\usepackage{amsmath, amssymb}
\usepackage{booktabs}
\usepackage{array}

% ── bảng màu ────────────────────────────────────────────────────────
\definecolor{navy}{HTML}{0E3F92}
\definecolor{tint}{HTML}{EDF2FB}
\definecolor{ink}{HTML}{2B2B2B}
\definecolor{mute}{HTML}{6B7280}
\definecolor{warn}{HTML}{9A3412}
\definecolor{good}{HTML}{14532D}

% ── theme tối giản ─────────────────────────────────────────────────
\usetheme{default}
\usecolortheme{default}
\setbeamertemplate{navigation symbols}{}
\setbeamercolor{normal text}{fg=ink}
\setbeamercolor{structure}{fg=navy}
\setbeamercolor{frametitle}{fg=white, bg=navy}
\setbeamerfont{frametitle}{size=\large, series=\bfseries}
\setbeamerfont{framesubtitle}{size=\small, series=\mdseries}
\setbeamercolor{framesubtitle}{fg=mute}

\setbeamertemplate{frametitle}{%
  \nointerlineskip
  \begin{beamercolorbox}[wd=\paperwidth, ht=2.4ex, dp=1.2ex, leftskip=0.9em]{frametitle}
    \usebeamerfont{frametitle}\insertframetitle
  \end{beamercolorbox}}

\setbeamertemplate{footline}{%
  \hbox{\begin{beamercolorbox}[wd=\paperwidth, ht=2.2ex, dp=1ex]{}
    \hspace{0.9em}\color{mute}\scriptsize\insertsection
    \hfill\color{mute}\scriptsize\insertframenumber\,/\,\inserttotalframenumber
    \hspace{0.9em}\end{beamercolorbox}}}

\setbeamertemplate{itemize item}{\color{navy}\small$\bullet$}
\setbeamertemplate{itemize subitem}{\color{mute}\scriptsize$\circ$}
\setlength{\leftmargini}{1.1em}
\setlength{\leftmarginii}{1.1em}

% ── hộp ─────────────────────────────────────────────────────────────
\tcbset{
  bx/.style={colback=tint, colframe=navy!22, boxrule=0.4pt, arc=2pt,
             left=4.5pt, right=4.5pt, top=2pt, bottom=2pt,
             fonttitle=\bfseries\small, coltitle=navy},
  key/.style={bx, colback=navy!6, colframe=navy!45},
  bad/.style={bx, colback=warn!6, colframe=warn!35, coltitle=warn},
}
\newtcolorbox{bx}[1][]{bx,#1}
\newtcolorbox{keyb}[1][]{key,#1}
\newtcolorbox{badb}[1][]{bad,#1}

\DefineVerbatimEnvironment{dia}{Verbatim}%
  {fontsize=\scriptsize, baselinestretch=0.92}
\newcommand{\dd}[1]{{\color{mute}\scriptsize #1}}
\newcommand{\hl}[1]{{\color{navy}\bfseries #1}}
\newcommand{\no}[1]{{\color{warn}\bfseries #1}}
\newcommand{\yes}[1]{{\color{good}\bfseries #1}}

\title{\bfseries quy gán nguồn cho rag hai tầng}
\subtitle{tập đủ tối tiểu \;$+$\; integrated gradients trên đồ thị}
\date{}
\author{}

\begin{document}

% ════════════════════════════════════════════════════════════════════
\begin{frame}[plain]
  \vfill
  {\color{navy}\Large\bfseries quy gán nguồn cho rag hai tầng}\\[0.9em]
  {\large tập đủ tối tiểu \;$+$\; integrated gradients trên đồ thị}\\[2.2em]
  \begin{tabular}{@{}ll@{}}
    \color{mute}tầng a & tìm \hl{tập đủ tối tiểu} trong pool ứng viên\\
                       & \dd{mwc · deegan–packel · banzhaf có dấu}\\[0.5em]
    \color{mute}tầng b & tìm \hl{đường} trên đồ thị tri thức\\
                       & \dd{integrated gradients trên gate của cạnh}\\
  \end{tabular}
  \vfill
\end{frame}

% ════════════════════════════════════════════════════════════════════
\section{bài toán}

\begin{frame}{bài toán}
  rag trả lời đúng, nhưng không nói được \emph{vì sao}. hai câu hỏi khác nhau:

  \vspace{0.6em}
  \begin{columns}[t]
  \begin{column}{0.48\textwidth}
    \begin{bx}[title={đoạn nào gánh câu trả lời}]
      trong $n$ đoạn đã lấy về, tập nào \emph{thật sự} đủ để trả lời,
      và không đoạn nào thừa?
      \vspace{0.3em}
      \dd{prompt phình theo số đoạn $\Rightarrow$ ablation}
    \end{bx}
  \end{column}
  \begin{column}{0.48\textwidth}
    \begin{bx}[title={vì sao truy hồi tìm ra chúng}]
      trên đồ thị tri thức, \emph{đường} nào dẫn từ thực thể trong câu hỏi
      tới đoạn đó?
      \vspace{0.3em}
      \dd{đồ thị cố định, hàng nghìn cạnh $\Rightarrow$ ig}
    \end{bx}
  \end{column}
  \end{columns}

  \vspace{0.8em}
  \begin{keyb}
    hai câu hỏi cần hai công cụ khác nhau. ghép chúng bằng \hl{một giao diện duy nhất}:
    tầng a trả về tập đích $T$, tầng b nhận $T$. \no{không điểm số nào đi qua.}
  \end{keyb}
\end{frame}

\begin{frame}{toàn cảnh}
\begin{dia}
truy vấn q
   |
   |  3 nguồn -> pool thô K <= 60 -> khử trùng lặp -> rerank THEO TỪNG NGUỒN
   v
pool C(q),  n = 40  =  20 bm25  +  15 gfm  +  5 ngẫu nhiên
   |
   |  >>> TẦNG A   tìm tập đủ tối tiểu
   |      v(S)    model đọc S rồi trả lời — khớp đáp án vàng không?
   |      adq(S)  các đoạn trong S có THẬT SỰ nói ra đáp án không?
   |              (adq dùng bảng phủ, 0 lượt LLM)
   v
T = S*  =  {đoạn 3, đoạn 12}          <- THỨ DUY NHẤT đi qua ranh giới
   |
   |  >>> TẦNG B   tìm đường trên đồ thị
   |      gate trên cạnh · ig trên gate · quy hoạch động · trích đường
   v
Pi_D  =  {đường 1, đường 2}  +  suff(Pi_D)  +  credit từng đoạn
\end{dia}
\end{frame}

